% Główny plik dokumentu – kompiluj właśnie ten plik

\documentclass[a4paper,12pt,oneside,onecolumn]{book}
% Ustawienia formatu dokumentu:
% - a4paper: format A4
% - 12pt: wielkość czcionki
% - oneside: jednostronny druk
% - onecolumn: jedna kolumna tekstu

\linespread{1.5} % Zwiększenie interlinii (1.5x) dla lepszej czytelności
\input{style.tex} % Dołączenie pliku z definicją stylu dokumentu

% Dołączenie plików z bibliografią (literatura i netografia)
\addbibresource{refs.bib}
\addbibresource{url_refs.bib}

\usepackage{adjustbox}
\usepackage{float}
\usepackage[titles]{tocloft}


\begin{document}

% *************** Front matter ***************
% *************** Strona tytułowa ***************

\pagestyle{empty}
\noindent
\begin{center}
    UNIWERSYTET EKONOMICZNY W KATOWICACH
\end{center}

\begin{center}
    INFORMATYKA
\end{center}

\vfill
\begin{center}
    \textbf{Patryk Klimek,}
    \textbf{Dominik Rekowski,}
    \textbf{Łukasz Osadnik}
\end{center}
\vfill

\begin{center}
    \Large
	\textbf{Rozkład jazdy autobusów dla miasta Sławków}
\end{center}

\begin{center}
	\textbf{Bus timetable for the city of Sławków}
\end{center}

\vfill
\vfill
\vfill


\vfill
\begin{center}
    KATOWICE 2026
\end{center}

\newpage 

% *************** Spis treści ***************
\tableofcontents

% *************** Koniec front matter *************** % Tutaj znajduje się strona tytułowa, spis treści, oświadczenia itd.

% *************** Treść rozdziałów *******-> 80-130 tys znakow

% \chapter*{Wstęp}
% -*- TeX-master: "../main.tex" -*-
\chapter{Wstęp do tematu}

Sławków to miasto, któremu brak autobusów. To sie zmienia z dniem dzisiejszym.
% -*- TeX-master: "../main.tex" -*-
\chapter{Wybór modelu autobusu}
Zaprojektowanie sieci przystanków i rozkładów jazdy musi iść w parze z odpowiednim doborem taboru, który fizycznie obsłuży zaplanowane trasy. Wybór modelu autobusu jest złożonym problemem decyzyjnym, wymagającym uwzględnienia ograniczeń budżetowych, wymagań terenowych oraz oczekiwań pasażerów. Ponieważ żaden dostępny na rynku pojazd nie jest idealny pod każdym względem (np. pojazdy najtańsze nie zawsze oferują najwyższy prestiż lub odpowiednią moc), konieczne było zastosowanie narzędzi z zakresu badań operacyjnych.

W niniejszym rozdziale przeprowadzono ocenę dostępnych na rynku autobusów z wykorzystaniem metody SAW (\textit{Simple Additive Weighting}). Pozwala ona na sprowadzenie różnorodnych parametrów technicznych i ekonomicznych do wspólnego mianownika, ułatwiając tym samym obiektywny wybór najlepszego rozwiązania dla Sławkowa.

\section{Założenia metody i wagi kryteriów}
Wychodząc naprzeciw specyficznym potrzebom nowo projektowanego układu komunikacyjnego w Sławkowie, parametry modelu decyzyjnego dostosowano w sposób odzwierciedlający realne wyzwania topograficzne oraz politykę transportową miasta. W procesie zrezygnowano z równego podziału wag, silnie premiując cechy gwarantujące sprawność obsługi i zadowolenie podróżnych. Wyodrębniono następujące kryteria:

\begin{itemize}
    \item \textbf{K1: Moc silnika [KM] (stymulanta, waga $0,26$)} -- podwyższona waga tego kryterium wynika bezpośrednio z topografii miasta Sławków. Nierówny teren i przewyższenia sprawiają, że odpowiednia moc napędowa jest absolutnie kluczowa do sprawnego pokonywania wzniesień i utrzymania płynności rozkładu jazdy.
    \item \textbf{K2: Miejsca dla pasażerów (stymulanta, waga $0,26$)} -- wysoka waga ma na celu zminimalizowanie zjawiska przepełnienia i "ciśnienia się" pasażerów w godzinach szczytowych, co bezpośrednio przekłada się na chęć korzystania z transportu publicznego.
    \item \textbf{K3: Prestiż i komfort [skala 1-10] (stymulanta, waga $0,26$)} -- subiektywna ocena uśredniona przez grono decydentów w sekcji \ref{sec:bus_prestige}. Uwzględnia wygląd zewnętrzny, estetykę wnętrza i ciekawe rozwiązania technologiczne. Priorytetyzacja tego aspektu wynika z chęci zapewnienia mieszkańcom jak najlepszych doświadczeń z podróży.
    \item \textbf{K4: Udogodnienia dla niepełnosprawnych (stymulanta, waga $0,11$)} -- kryterium dyskretne zliczające łączną ilość zastosowanych w pojeździe ułatwień dla osób o ograniczonej mobilności (np. rampa, niska podłoga, przyklęk, miejsce na wózek, braille).
    \item \textbf{K5: Koszt zakupu [mln PLN] (destymulanta, waga $0,11$)} -- szacunkowy koszt jednostkowy pojazdu.
\end{itemize}

Aby zapewnić miarodajność, analizę rozdzielono na dwie niezależne kategorie rynkowe: autobusy duże (przegubowe i wielkopojemne) oraz autobusy średnie (pojazdy klasy MIDI).

\section{Ocena prestiżu autobusów}\label{sec:bus_prestige}

Jednym z kluczowych kryteriów przyjętych w analizie decyzyjnej jest prestiż pojazdu. Wartość tego parametru wyznaczono jako średnią arytmetyczną z ocen eksperckich zespołu projektowego. Składowe oceny obejmowały analizę designu zewnętrznego, aranżacji wnętrza, a także identyfikację niestandardowych udogodnień i nowatorskich rozwiązań konstrukcyjnych.

Dodatkowym atutem, który bezpośrednio wpłynął na podwyższenie punktacji dla pojazdów marki Solaris, jest obecność charakterystycznej maskotki jamnika. Zespół decyzyjny docenił ten nietypowy element wizerunkowy, ponieważ pozytywnie wpływa on na odbiór estetyczny autobusu i kreuje przyjazną atmosferę dla podróżnych, szczególnie tych najmłodszych.
\MyFigure[0.4\textwidth]{images/solaris_jamnik.jpg}{Maskotka firmy Solaris}{fig:solarisjamnik}

Poniżej przedstawiono charakterystykę wszystkich wytypowanych modeli autobusów wraz z merytorycznym uzasadnieniem przyznanych im not.

\subsection{Solaris Urbino 15 LE electric}
Decydenci wysoko ocenili nowoczesny design zewnętrza pojazdu, oraz minimalistyczny projekt przestrzeni pasażerskiej. Znaczącym atutem pojazdu w kategorii komfortu są ergonomicznie wyprofilowane fotele, oferujące stabilne podparcie kręgosłupa podczas podróży.
\MyFigure[0.5\textwidth]{images/autobusy/Urbino15LEelectric/wnetrze.jpg}{Wnętrze autobusu Urbino 15 LE electric}{fig:Urbino15LEelectricwnetrze}
\MyFigure[0.5\textwidth]{images/autobusy/Urbino15LEelectric/zewnatrz.png}{Zewnętrze autobusu Urbino 15 LE electric}{fig:Urbino15LEelectriczewnatrz}

\subsection{CapaCity L}
Wysoka nota wynika przede wszystkim z doskonale zaprojektowanego wnętrza. Decydenci zwrócili szczególną uwagę na przestronność oraz imponującą liczbę dostępnych foteli. Całokształt aranżacji sprawia, że standard przejazdu przypomina podróż komfortowym autokarem wycieczkowym, co bezpośrednio przełoży się na wysoką satysfakcję pasażerów.
\MyFigure[0.5\textwidth]{images/autobusy/CapaCityL/wnetrze.png}{Wnętrze autobusu CapaCity L}{fig:CapaCityLwnetrze}
\MyFigure[0.5\textwidth]{images/autobusy/CapaCityL/zewnatrz.png}{Zewnętrze autobusu CapaCity L}{fig:CapaCityLzewnatrz}

\subsection{Citaro K}
Wysoka nota przyznana temu pojazdowi wynika w dużej mierze z bardzo udanego designu nadwozia. Front autobusu odznacza się masywną, solidną sylwetką, która budzi poczucie trwałości i bezpieczeństwa. Linia boczna harmonijnie współgra z estetyczną kolorystyką powłoki lakierniczej (w wariancie poglądowym jest to elegancki, złoty odcień). Kolejnym kluczowym atutem jest przemyślana aranżacja wnętrza. Zespół decyzyjny docenił układ tylnych rzędów siedzeń, których oparcia posiadają lekkie pochylenie, gwarantując pasażerom wyższy poziom odprężenia i komfortu w trakcie podróży.
\MyFigure[0.5\textwidth]{images/autobusy/CitaroK/wnetrze.png}{Wnętrze autobusu Citaro K}{fig:CitaroKwnetrze}
\MyFigure[0.5\textwidth]{images/autobusy/CitaroK/zewnatrz.png}{Zewnętrze autobusu Citaro K}{fig:CitaroKzewnatrz}

\subsection{Crossway 10.8 LE}
Niezwykle wysoka nota przyznana temu pojazdowi jest wynikiem jego wyjątkowych walorów estetycznych i funkcjonalnych. Zespół decyzyjny docenił przede wszystkim nowoczesną, modernistyczną linię nadwozia. Przestrzeń pasażerska wyróżnia się wysoką spójnością kolorystyczną oraz ogólną elegancją aranżacji. Na szczególne uznanie zasługują praktyczne udogodnienia w postaci dedykowanych przestrzeni bagażowych nad siedzeniami. Ponadto, zastosowanie podwyższonego podestu w tylnej sekcji pasażerskiej nadaje tej strefie ekskluzywny charakter, co znacząco podnosi prestiż i odczuwalny komfort podróży.
\MyFigure[0.5\textwidth]{images/autobusy/Crossway10.8LE/wnetrze.png}{Wnętrze autobusu Crossway 10.8 LE}{fig:Crossway108LEwnetrze}
\MyFigure[0.5\textwidth]{images/autobusy/Crossway10.8LE/zewnatrz.png}{Zewnętrze autobusu Crossway 10.8 LE}{fig:Crossway108LEzewnatrz}

\subsection{Crossway Pro 10,8}
Pojazd ten został oceniony na 6 punktów ze względu na swój surowy charakter. Decydenci określili design zewnętrzny jako nijaki i przestarzały, co w zestawieniu z proporcjami okien nadaje autobusowi ciężki, wizualnie zamknięty wygląd. Wnętrze pojazdu kontynuuje ten sam trend - aranżacja siedzeń, ułożonych w ścisłych, wysokich rzędach, sprawia wrażenie przytłaczającej i brakuje w niej elementów ocieplających wizerunek. W efekcie podróżni mogą odczuwać brak przestrzeni i swobody.
\MyFigure[0.5\textwidth]{images/autobusy/CrosswayPro10,8/wnetrze.png}{Wnętrze autobusu Crossway Pro 10,8}{fig:CrosswayPro108wnetrze}
\MyFigure[0.5\textwidth]{images/autobusy/CrosswayPro10,8/zewnatrz.png}{Zewnętrze autobusu Crossway Pro 10,8}{fig:CrosswayPro108zewnatrz}

\subsection{MAN Lion's City 10 E}
Pojazd ten otrzymał niższą punktację głównie ze względu na dysonans między poprawnym designem nadwozia a dyskusyjnym projektem przestrzeni pasażerskiej. Decydenci zwrócili szczególną uwagę na zbyt wąskie siedziska i oparcia foteli, które nie gwarantują odpowiedniego komfortu i mogą powodować uciążliwy ścisk w ułożeniu dwurzędowym. Dodatkowym czynnikiem negatywnie rzutującym na estetykę wnętrza jest nietypowe wykończenie podłogi. Wzór przypominający tradycyjne, domowe panele podłogowe został oceniony jako rozwiązanie nienaturalne i zaburzające profesjonalny wizerunek transportu zbiorowego.
\MyFigure[0.5\textwidth]{images/autobusy/MANLion'sCity10E/wnetrze.png}{Wnętrze autobusu MAN Lion's City 10 E}{fig:MANLionsCity10Ewnetrze}
\MyFigure[0.5\textwidth]{images/autobusy/MANLion'sCity10E/zewnatrz.png}{Zewnętrze autobusu MAN Lion's City 10 E}{fig:MANLionsCity10Ezewnatrz}

\subsection{MAN Lion's City 12 E}
Model ten zachowuje atrakcyjnie zaprojektowaną bryłę zewnętrzną, docenioną już w pokrewnych pojazdach tego producenta. Jego zdecydowaną przewagą jest natomiast znacznie przestronniejsze i bardziej funkcjonalne wnętrze. Zespół decyzyjny zauważył, że dodatkowa przestrzeń w kabinie pasażerskiej skutecznie poprawia ergonomię, oferując podróżnym o wiele wyższy poziom swobody oraz lepsze wrażenia z jazdy.
\MyFigure[0.5\textwidth]{images/autobusy/MANLion'sCity12E/wnetrze.png}{Wnętrze autobusu MAN Lion's City 12 E}{fig:MANLionsCity12Ewnetrze}
\MyFigure[0.5\textwidth]{images/autobusy/MANLion'sCity12E/zewnatrz.png}{Zewnętrze autobusu MAN Lion's City 12 E}{fig:MANLionsCity12Ezewnatrz}

\subsection{MAN Lion's City 12 G}
Niska ocena prestiżu jest w tym przypadku konsekwencją przeciętnej estetyki nadwozia oraz istotnych braków w wyposażeniu. Design zewnętrzny, choć w pełni akceptowalny, opiera się na powszechnie stosowanych schematach. Głównym powodem krytycznej oceny decydentów jest jednak wnętrze pojazdu. Kabina pasażerska charakteryzuje się niekorzystnym, surowym układem przestrzennym. Zainstalowane siedzenia wyglądają na niedostatecznie wyprofilowane, co negatywnie rzutuje na postrzegany komfort podróży. Dodatkowo, dużą wadą w tej kategorii okazał się brak zintegrowanych paneli wizualnej informacji dla podróżnych, które stanowią obecnie bazowy standard rynkowy.
\MyFigure[0.5\textwidth]{images/autobusy/MANLion'sCity12G/wnetrze.png}{Wnętrze autobusu MAN Lion's City 12 G}{fig:MANLionsCity12Gwnetrze}
\MyFigure[0.5\textwidth]{images/autobusy/MANLion'sCity12G/zewnatrz.png}{Zewnętrze autobusu MAN Lion's City 12 G}{fig:MANLionsCity12Gzewnatrz}

\subsection{MAN Lion's Intercity LE 14}
Model ten uzyskał umiarkowaną ocenę na poziomie 6 punktów. Choć ogólna linia nadwozia jest w dużej mierze poprawna, zespół decyzyjny krytycznie ocenił charakterystyczny, wydłużony pas przedni pojazdu, który wizualnie zaburza proporcje bryły zewnętrznej. Wnętrze autobusu wzbudziło mieszane odczucia. Z jednej strony pozytywnie oceniono podwyższoną sekcję tylną, urozmaicającą przestrzeń pasażerską. Z drugiej jednak, ocenę obniżyło ponowne zastosowanie dyskusyjnego wykończenia podłogi imitującego domowe panele drewniane. Istotnym zarzutem z perspektywy ergonomii jest również zbyt gęsty rozstaw foteli. Zredukowana odległość między rzędami znacząco ogranicza przestrzeń na nogi, co może powodować dyskomfort, w szczególności u osób o wyższym wzroście.
\MyFigure[0.4\textwidth]{images/autobusy/MANLion'sIntercityLE14/wnetrze.png}{Wnętrze autobusu MAN Lion's Intercity LE 14}{fig:MANLionsIntercityLE14wnetrze}
\MyFigure[0.5\textwidth]{images/autobusy/MANLion'sIntercityLE14/zewnatrz.png}{Zewnętrze autobusu MAN Lion's Intercity LE 14}{fig:MANLionsIntercityLE14zewnatrz}

\subsection{Nesobus 12}
Prezentowany model to niekwestionowany faworyt zespołu decyzyjnego w kategorii prestiżu. Zewnętrzna bryła pojazdu przyciąga uwagę modernistycznym, niezwykle nowoczesnym designem. Kluczowym czynnikiem, który zadecydował o tak wysokiej nocie, jest jednak bezprecedensowa jakość aranżacji wnętrza. Przestrzeń pasażerska odznacza się elegancją oraz harmonijnie dobraną paletą barw. Co istotne, zastosowana tu wykładzina drewnopodobna – w przeciwieństwie do innych ocenianych modeli – nie sprawia wrażenia taniego wykończenia. Wraz z architekturą siedzeń kreuje ona wysoce prestiżową, wręcz luksusową atmosferę. Na szczególne uznanie zasługuje innowacyjne zagospodarowanie strefy dla pasażerów stojących, obejmujące rozkładane fotele oraz tapicerowane, miękkie oparcia. Całości dopełnia wysoce ergonomiczny i optymalny układ standardowych miejsc siedzących.
\MyFigure[0.5\textwidth]{images/autobusy/Nesobus12/wnetrze.png}{Wnętrze autobusu Nesobus 12}{fig:Nesobus12wnetrze}
\MyFigure[0.5\textwidth]{images/autobusy/Nesobus12/zewnatrz.jpg}{Zewnętrze autobusu Nesobus 12}{fig:Nesobus12zewnatrz}

\subsection{SANCITY 9LE}
Model ten uzyskał od zespołu decyzyjnego stosunkowo niską ocenę na poziomie 4 punktów. O ile stylistyka zewnętrzna pojazdu prezentuje się estetycznie i nie budzi większych zastrzeżeń, o tyle wnętrze spotkało się z surową krytyką. Przestrzeń pasażerska sprawia wrażenie monotonnej, wizualnie pustej i pozbawionej wyrazu. Zauważalna jest nieoptymalna aranżacja kabiny, która skutkuje wyraźnym i nieuzasadnionym deficytem miejsc siedzących względem dostępnej kubatury. Jedynym elementem architektonicznym, który w pewnym stopniu urozmaica to surowe wnętrze i zapobiega przyznaniu jeszcze niższej noty, jest zastosowanie podwyższenia w tylnej sekcji pojazdu.
\MyFigure[0.5\textwidth]{images/autobusy/SANCITY9LE/wnetrze.png}{Wnętrze autobusu SANCITY 9LE}{fig:SANCITY9LEwnetrze}
\MyFigure[0.5\textwidth]{images/autobusy/SANCITY9LE/zewnatrz.png}{Zewnętrze autobusu SANCITY 9LE}{fig:SANCITY9LEzewnatrz}

\subsection{SANCITY 10LF}
Miejsce na opis i uzasadnienie oceny prestiżu dla tego modelu...
\MyFigure[0.5\textwidth]{images/autobusy/SANCITY10LF/wnetrze.png}{Wnętrze autobusu SANCITY 10LF}{fig:SANCITY10LFwnetrze}
\MyFigure[0.5\textwidth]{images/autobusy/SANCITY10LF/zewnatrz.png}{Zewnętrze autobusu SANCITY 10LF}{fig:SANCITY10LFzewnatrz}

\subsection{SANCITY 10LFE}
Miejsce na opis i uzasadnienie oceny prestiżu dla tego modelu...
\MyFigure[0.4\textwidth]{images/autobusy/SANCITY10LFE/wnetrze.png}{Wnętrze autobusu SANCITY 10LFE}{fig:SANCITY10LFEwnetrze}
\MyFigure[0.5\textwidth]{images/autobusy/SANCITY10LFE/zewnatrz.png}{Zewnętrze autobusu SANCITY 10LFE}{fig:SANCITY10LFEzewnatrz}

\subsection{SANCITY 12LF}
Miejsce na opis i uzasadnienie oceny prestiżu dla tego modelu...
\MyFigure[0.5\textwidth]{images/autobusy/SANCITY12LF/wnetrze.png}{Wnętrze autobusu SANCITY 12LF}{fig:SANCITY12LFwnetrze}
\MyFigure[0.5\textwidth]{images/autobusy/SANCITY12LF/zewnatrz.png}{Zewnętrze autobusu SANCITY 12LF}{fig:SANCITY12LFzewnatrz}

\subsection{SANCITY 12LFH}
Miejsce na opis i uzasadnienie oceny prestiżu dla tego modelu...
\MyFigure[0.5\textwidth]{images/autobusy/SANCITY12LFH/wnetrze.png}{Wnętrze autobusu SANCITY 12LFH}{fig:SANCITY12LFHwnetrze}
\MyFigure[0.5\textwidth]{images/autobusy/SANCITY12LFH/zewnatrz.png}{Zewnętrze autobusu SANCITY 12LFH}{fig:SANCITY12LFHzewnatrz}

\subsection{Urbino 8,9 LE}
Miejsce na opis i uzasadnienie oceny prestiżu dla tego modelu...
\MyFigure[0.5\textwidth]{images/autobusy/Urbino8,9LE/wnetrze.png}{Wnętrze autobusu Urbino 8,9 LE}{fig:Urbino89LEwnetrze}
\MyFigure[0.5\textwidth]{images/autobusy/Urbino8,9LE/zewnatrz.png}{Zewnętrze autobusu Urbino 8,9 LE}{fig:Urbino89LEzewnatrz}

\subsection{Urbino 8,9 LE electric}
Miejsce na opis i uzasadnienie oceny prestiżu dla tego modelu...
\MyFigure[0.5\textwidth]{images/autobusy/Urbino8,9LEelectric/wnetrze.png}{Wnętrze autobusu Urbino 8,9 LE electric}{fig:Urbino89LEelectricwnetrze}
\MyFigure[0.5\textwidth]{images/autobusy/Urbino8,9LEelectric/zewnatrz.png}{Zewnętrze autobusu Urbino 8,9 LE electric}{fig:Urbino89LEelectriczewnatrz}

\subsection{Urbino 10,5 electric}
Miejsce na opis i uzasadnienie oceny prestiżu dla tego modelu...
\MyFigure[0.5\textwidth]{images/autobusy/Urbino10,5electric/wnetrze.png}{Wnętrze autobusu Urbino 10,5 electric}{fig:Urbino105electricwnetrze}
\MyFigure[0.5\textwidth]{images/autobusy/Urbino10,5electric/zewnatrz.jpg}{Zewnętrze autobusu Urbino 10,5 electric}{fig:Urbino105electriczewnatrz}

\subsection{Urbino 12}
Miejsce na opis i uzasadnienie oceny prestiżu dla tego modelu...
\MyFigure[0.5\textwidth]{images/autobusy/Urbino12/wnetrze.png}{Wnętrze autobusu Urbino 12}{fig:Urbino12wnetrze}
\MyFigure[0.5\textwidth]{images/autobusy/Urbino12/zewnatrz.jpg}{Zewnętrze autobusu Urbino 12}{fig:Urbino12zewnatrz}

\subsection{Urbino 12 CNG}
Miejsce na opis i uzasadnienie oceny prestiżu dla tego modelu...
\MyFigure[0.5\textwidth]{images/autobusy/Urbino12CNG/wnetrze.png}{Wnętrze autobusu Urbino 12 CNG}{fig:Urbino12CNGwnetrze}
\MyFigure[0.5\textwidth]{images/autobusy/Urbino12CNG/zewnatrz.png}{Zewnętrze autobusu Urbino 12 CNG}{fig:Urbino12CNGzewnatrz}

\subsection{Urbino 15 LE electric}
Miejsce na opis i uzasadnienie oceny prestiżu dla tego modelu...
\MyFigure[0.5\textwidth]{images/autobusy/Urbino15LEelectric/wnetrze.jpg}{Wnętrze autobusu Urbino 15 LE electric}{fig:Urbino15LEelectricwnetrze}
\MyFigure[0.5\textwidth]{images/autobusy/Urbino15LEelectric/zewnatrz.png}{Zewnętrze autobusu Urbino 15 LE electric}{fig:Urbino15LEelectriczewnatrz}

\subsection{Urbino 18}
Miejsce na opis i uzasadnienie oceny prestiżu dla tego modelu...
\MyFigure[0.5\textwidth]{images/autobusy/Urbino18/wnetrze.png}{Wnętrze autobusu Urbino 18}{fig:Urbino18wnetrze}
\MyFigure[0.5\textwidth]{images/autobusy/Urbino18/zewnatrz.jpg}{Zewnętrze autobusu Urbino 18}{fig:Urbino18zewnatrz}

\subsection{Urbino 18 CNG}
Miejsce na opis i uzasadnienie oceny prestiżu dla tego modelu...
\MyFigure[0.5\textwidth]{images/autobusy/Urbino18CNG/wnetrze.png}{Wnętrze autobusu Urbino 18 CNG}{fig:Urbino18CNGwnetrze}
\MyFigure[0.7\textwidth]{images/autobusy/Urbino18CNG/zewnatrz.png}{Zewnętrze autobusu Urbino 18 CNG}{fig:Urbino18CNGzewnatrz}

\section{Analiza SAW dla autobusów dużych}

Do tej kategorii zakwalifikowano pojazdy o wysokiej zdolności przewozowej, przeznaczone na główne osie komunikacyjne miasta.

\begin{table}[H]
    \centering
    \caption{Macierz decyzyjna (dane wejściowe) -- autobusy duże}
    \vspace{0.2cm}
    \begin{adjustbox}{max width=\textwidth}
    \begin{tabular}{|l|c|c|c|c|c|}
        \hline
        \textbf{Model Autobusu} & \textbf{K1: Moc [KM]} & \textbf{K2: Miejsca} & \textbf{K3: Prestiż} & \textbf{K4: Udogod.} & \textbf{K5: Koszt [mln]} \\ 
        & \textbf{(MAX)} & \textbf{(MAX)} & \textbf{(MAX)} & \textbf{(MAX)} & \textbf{(MIN)} \\ \hline
        Urbino 15 LE electric & 408 & 65 & 7.00 & 2 & 2.20 \\ \hline
        Urbino 18 CNG & 320 & 55 & 10.00 & 3 & 2.40 \\ \hline
        Urbino 18 & 310 & 54 & 6.33 & 3 & 1.80 \\ \hline
        Urbino 12 & 256 & 16 & 8.67 & 3 & 1.00 \\ \hline
        Urbino 12 CNG & 320 & 40 & 4.67 & 3 & 1.10 \\ \hline
        MAN Lion's City 12 E & 326 & 44 & 6.33 & 4 & 2.60 \\ \hline
        MAN Lion's Intercity LE 14 & 265 & 63 & 6.33 & 3 & 1.80 \\ \hline
        MAN Lion's City 12 G & 280 & 37 & 3.00 & 4 & 2.00 \\ \hline
        SANCITY 12LFH & 344 & 40 & 8.00 & 3 & 3.25 \\ \hline
        SANCITY 12LF & 277 & 40 & 4.00 & 3 & 1.00 \\ \hline
        CapaCity L & 360 & 44 & 8.67 & 4 & 2.00 \\ \hline
        Nesobus 12 & 340 & 37 & 9.67 & 3 & 2.90 \\ \hline
    \end{tabular}
    \end{adjustbox}
\end{table}

\begin{table}[H]
    \centering
    \caption{Wyniki modelu SAW oraz ranking -- autobusy duże}
    \vspace{0.2cm}
    \begin{adjustbox}{max width=\textwidth}
    \begin{tabular}{|c|l|c|c|c|c|c|c|}
        \hline
        \textbf{Poz.} & \textbf{Model Autobusu} & \textbf{K1 norm.} & \textbf{K2 norm.} & \textbf{K3 norm.} & \textbf{K4 norm.} & \textbf{K5 norm.} & \textbf{Wynik SAW} \\ \hline
        1 & Urbino 18 CNG & 0.78 & 0.85 & 1.00 & 0.75 & 0.42 & \textbf{0.812} \\ \hline
        2 & Urbino 15 LE electric & 1.00 & 1.00 & 0.70 & 0.50 & 0.45 & \textbf{0.807} \\ \hline
        3 & CapaCity L & 0.88 & 0.68 & 0.87 & 1.00 & 0.50 & \textbf{0.796} \\ \hline
        4 & Nesobus 12 & 0.83 & 0.57 & 0.97 & 0.75 & 0.34 & \textbf{0.736} \\ \hline
        5 & MAN Lion's Intercity LE 14 & 0.65 & 0.97 & 0.63 & 0.75 & 0.56 & \textbf{0.729} \\ \hline
        6 & Urbino 18 & 0.76 & 0.83 & 0.63 & 0.75 & 0.56 & \textbf{0.722} \\ \hline
        7 & SANCITY 12LFH & 0.84 & 0.62 & 0.80 & 0.75 & 0.31 & \textbf{0.704} \\ \hline
        8 & MAN Lion's City 12 E & 0.80 & 0.68 & 0.63 & 1.00 & 0.38 & \textbf{0.701} \\ \hline
        9 & Urbino 12 CNG & 0.78 & 0.62 & 0.47 & 0.75 & 0.91 & \textbf{0.668} \\ \hline
        10 & Urbino 12 & 0.63 & 0.25 & 0.87 & 0.75 & 1.00 & \textbf{0.645} \\ \hline
        11 & SANCITY 12LF & 0.68 & 0.62 & 0.40 & 0.75 & 1.00 & \textbf{0.633} \\ \hline
        12 & MAN Lion's City 12 G & 0.69 & 0.57 & 0.30 & 1.00 & 0.50 & \textbf{0.569} \\ \hline
    \end{tabular}
    \end{adjustbox}
\end{table}

Skupienie wagi na prestiżu i ergonomii pasażerskiej doprowadziło do zwycięstwa pojazdu z napędem gazowym -- \textit{Solaris Urbino 18 CNG}. Zgromadził on doskonałe noty za estetykę i oferowane rozwiązania (prestiż uśredniony na 10,00) oraz dużą pojemność. Bezpośrednio za nim znalazł się najpotężniejszy (408 KM) elektryczny \textit{Urbino 15 LE}.

\section{Analiza SAW dla autobusów średnich}

Zestawienie pojazdów o długości do kilkunastu metrów, które idealnie wpisują się w linie dowozowe, osiedlowe oraz trasy o węższej infrastrukturze drogowej.

\begin{table}[H]
    \centering
    \caption{Macierz decyzyjna (dane wejściowe) -- autobusy średnie}
    \vspace{0.2cm}
    \begin{adjustbox}{max width=\textwidth}
    \begin{tabular}{|l|c|c|c|c|c|}
        \hline
        \textbf{Model Autobusu} & \textbf{K1: Moc [KM]} & \textbf{K2: Miejsca} & \textbf{K3: Prestiż} & \textbf{K4: Udogod.} & \textbf{K5: Koszt [mln]} \\ 
        & \textbf{(MAX)} & \textbf{(MAX)} & \textbf{(MAX)} & \textbf{(MAX)} & \textbf{(MIN)} \\ \hline
        Urbino 10,5 electric & 340 & 33 & 4.00 & 3 & 2.60 \\ \hline
        Urbino 8,9 LE electric & 218 & 27 & 8.00 & 4 & 2.40 \\ \hline
        Urbino 8,9 LE & 250 & 33 & 6.67 & 4 & 1.10 \\ \hline
        MAN Lion's City 10 E & 326 & 36 & 5.00 & 4 & 2.60 \\ \hline
        SANCITY 9LE & 210 & 16 & 4.33 & 4 & 0.60 \\ \hline
        SANCITY 10LFE & 340 & 24 & 5.00 & 4 & 2.20 \\ \hline
        SANCITY 10LF & 247 & 26 & 5.67 & 3 & 0.90 \\ \hline
        Citaro K & 299 & 32 & 8.33 & 5 & 1.30 \\ \hline
        Crossway 10.8LE & 330 & 29 & 8.67 & 4 & 1.10 \\ \hline
        Crossway Pro 10,8 & 360 & 47 & 6.33 & 3 & 0.90 \\ \hline
    \end{tabular}
    \end{adjustbox}
\end{table}

\begin{table}[H]
    \centering
    \caption{Wyniki modelu SAW oraz ranking -- autobusy średnie}
    \vspace{0.2cm}
    \begin{adjustbox}{max width=\textwidth}
    \begin{tabular}{|c|l|c|c|c|c|c|c|}
        \hline
        \textbf{Poz.} & \textbf{Model Autobusu} & \textbf{K1 norm.} & \textbf{K2 norm.} & \textbf{K3 norm.} & \textbf{K4 norm.} & \textbf{K5 norm.} & \textbf{Wynik SAW} \\ \hline
        1 & Crossway Pro 10,8 & 1.00 & 1.00 & 0.73 & 0.60 & 0.67 & \textbf{0.849} \\ \hline
        2 & Crossway 10.8LE & 0.92 & 0.62 & 1.00 & 0.80 & 0.55 & \textbf{0.807} \\ \hline
        3 & Citaro K & 0.83 & 0.68 & 0.96 & 1.00 & 0.46 & \textbf{0.804} \\ \hline
        4 & Urbino 8,9 LE & 0.69 & 0.70 & 0.77 & 0.80 & 0.55 & \textbf{0.711} \\ \hline
        5 & MAN Lion's City 10 E & 0.91 & 0.77 & 0.58 & 0.80 & 0.23 & \textbf{0.698} \\ \hline
        6 & Urbino 8,9 LE electric & 0.61 & 0.57 & 0.92 & 0.80 & 0.25 & \textbf{0.662} \\ \hline
        7 & SANCITY 10LFE & 0.94 & 0.51 & 0.58 & 0.80 & 0.27 & \textbf{0.646} \\ \hline
        8 & Urbino 10,5 electric & 0.94 & 0.70 & 0.46 & 0.60 & 0.23 & \textbf{0.639} \\ \hline
        9 & SANCITY 10LF & 0.69 & 0.55 & 0.65 & 0.60 & 0.67 & \textbf{0.632} \\ \hline
        10 & SANCITY 9LE & 0.58 & 0.34 & 0.50 & 0.80 & 1.00 & \textbf{0.568} \\ \hline
    \end{tabular}
    \end{adjustbox}
\end{table}

W tym segmencie dominację pokazała marka Iveco za sprawą modelu \textit{Crossway Pro 10,8}. Oferuje on absolutnie bezkonkurencyjną moc silnika w swojej klasie gabarytowej (360 KM) połączoną z maksymalizacją wykorzystania wnętrza pojazdu (najwyższa ocena w parametrze "Miejsca"). Świetnie wyważony pod kątem prestiżu w stosunku do mocy okazał się też \textit{Mercedes-Benz Citaro K}, zyskując najwyższą notę za udogodnienia dla niepełnosprawnych.
\chapter{Analiza wrażliwości}

Analiza wrażliwości pozwala ocenić, w jakim stopniu zmiana struktury wag kryteriów wpływa na końcowy ranking autobusów wyznaczony metodą SAW. Dla każdej kategorii pojazdów zdefiniowano trzy scenariusze preferencji decydenta: \textbf{scenariusz oszczędny} (dominująca waga kosztu zakupu), \textbf{scenariusz techniczny} (dominujące wagi mocy silnika i liczby miejsc) oraz \textbf{scenariusz neutralny} (równe wagi wszystkich kryteriów). Zestawienie wyników umożliwia identyfikację alternatyw stabilnych rankingowo oraz tych, których pozycja silnie zależy od przyjętych preferencji.


\section{Duże autobusy}

W segmencie dużych autobusów przeanalizowano wpływ trzech układów wag na wyniki modelu SAW. Poszczególne scenariusze różnią się zasadniczym przesunięciem punktu ciężkości oceny -- od priorytetu ekonomicznego po preferencje techniczne.

\begin{table}[H]
    \centering
    \caption{Scenariusze wag -- autobusy duże}
    \vspace{0.2cm}
    \begin{tabular}{|l|c|c|c|c|c|}
        \hline
        \textbf{Scenariusz} & \textbf{K1: Moc} & \textbf{K2: Miejsca} & \textbf{K3: Prestiż} & \textbf{K4: Udogod.} & \textbf{K5: Koszt} \\ \hline
        Oszczędny  & 0.14 & 0.14 & 0.14 & 0.08 & 0.50 \\ \hline
        Techniczny & 0.35 & 0.35 & 0.10 & 0.15 & 0.05 \\ \hline
        Neutralny  & 0.20 & 0.20 & 0.20 & 0.20 & 0.20 \\ \hline
    \end{tabular}
\end{table}

\begin{table}[H]
    \centering
    \caption{Wyniki analizy wrażliwości -- autobusy duże (scenariusz oszczędny)}
    \vspace{0.2cm}
    \begin{adjustbox}{max width=\textwidth}
    \begin{tabular}{|c|l|c|c|c|c|c|c|}
        \hline
        \textbf{Poz.} & \textbf{Model Autobusu} & \textbf{K1} & \textbf{K2} & \textbf{K3} & \textbf{K4} & \textbf{K5} & \textbf{Wynik SAW} \\ \hline
        1 & Urbino 12              & 0.088 & 0.034 & 0.121 & 0.060 & 0.500 & \textbf{0.804} \\ \hline
        2 & SANCITY 12LF           & 0.095 & 0.086 & 0.056 & 0.060 & 0.500 & \textbf{0.797} \\ \hline
        3 & Urbino 12 CNG          & 0.110 & 0.086 & 0.065 & 0.060 & 0.455 & \textbf{0.776} \\ \hline
        4 & CapaCity L             & 0.124 & 0.095 & 0.121 & 0.080 & 0.250 & \textbf{0.670} \\ \hline
        5 & MAN Lion's Intercity LE 14 & 0.091 & 0.136 & 0.089 & 0.060 & 0.278 & \textbf{0.653} \\ \hline
        6 & Urbino 18              & 0.106 & 0.116 & 0.089 & 0.060 & 0.278 & \textbf{0.649} \\ \hline
        7 & Urbino 15 LE electric  & 0.140 & 0.140 & 0.098 & 0.040 & 0.227 & \textbf{0.645} \\ \hline
        8 & Urbino 18 CNG          & 0.110 & 0.118 & 0.140 & 0.060 & 0.208 & \textbf{0.637} \\ \hline
        9 & MAN Lion's City 12 E   & 0.112 & 0.095 & 0.089 & 0.080 & 0.192 & \textbf{0.568} \\ \hline
        10 & Nesobus 12            & 0.117 & 0.080 & 0.135 & 0.060 & 0.172 & \textbf{0.564} \\ \hline
        11 & MAN Lion's City 12 G  & 0.096 & 0.080 & 0.042 & 0.080 & 0.250 & \textbf{0.548} \\ \hline
        12 & SANCITY 12LFH         & 0.118 & 0.086 & 0.112 & 0.060 & 0.154 & \textbf{0.530} \\ \hline
    \end{tabular}
    \end{adjustbox}
\end{table}

\begin{table}[H]
    \centering
    \caption{Wyniki analizy wrażliwości -- autobusy duże (scenariusz techniczny)}
    \vspace{0.2cm}
    \begin{adjustbox}{max width=\textwidth}
    \begin{tabular}{|c|l|c|c|c|c|c|c|}
        \hline
        \textbf{Poz.} & \textbf{Model Autobusu} & \textbf{K1} & \textbf{K2} & \textbf{K3} & \textbf{K4} & \textbf{K5} & \textbf{Wynik SAW} \\ \hline
        1 & Urbino 15 LE electric      & 0.350 & 0.350 & 0.070 & 0.075 & 0.023 & \textbf{0.868} \\ \hline
        2 & CapaCity L                 & 0.309 & 0.237 & 0.087 & 0.150 & 0.025 & \textbf{0.807} \\ \hline
        3 & Urbino 18 CNG              & 0.275 & 0.296 & 0.100 & 0.113 & 0.021 & \textbf{0.804} \\ \hline
        4 & MAN Lion's Intercity LE 14 & 0.227 & 0.339 & 0.063 & 0.113 & 0.028 & \textbf{0.770} \\ \hline
        5 & Urbino 18                  & 0.266 & 0.291 & 0.063 & 0.113 & 0.028 & \textbf{0.760} \\ \hline
        6 & MAN Lion's City 12 E       & 0.280 & 0.237 & 0.063 & 0.150 & 0.019 & \textbf{0.749} \\ \hline
        7 & SANCITY 12LFH              & 0.295 & 0.215 & 0.080 & 0.113 & 0.015 & \textbf{0.718} \\ \hline
        8 & Nesobus 12                 & 0.292 & 0.199 & 0.097 & 0.113 & 0.017 & \textbf{0.717} \\ \hline
        9 & Urbino 12 CNG              & 0.275 & 0.215 & 0.047 & 0.113 & 0.045 & \textbf{0.695} \\ \hline
        10 & SANCITY 12LF              & 0.238 & 0.215 & 0.040 & 0.113 & 0.050 & \textbf{0.656} \\ \hline
        11 & MAN Lion's City 12 G      & 0.240 & 0.199 & 0.030 & 0.150 & 0.025 & \textbf{0.644} \\ \hline
        12 & Urbino 12                 & 0.220 & 0.086 & 0.087 & 0.113 & 0.050 & \textbf{0.555} \\ \hline
    \end{tabular}
    \end{adjustbox}
\end{table}

\begin{table}[H]
    \centering
    \caption{Wyniki analizy wrażliwości -- autobusy duże (scenariusz neutralny)}
    \vspace{0.2cm}
    \begin{adjustbox}{max width=\textwidth}
    \begin{tabular}{|c|l|c|c|c|c|c|c|}
        \hline
        \textbf{Poz.} & \textbf{Model Autobusu} & \textbf{K1} & \textbf{K2} & \textbf{K3} & \textbf{K4} & \textbf{K5} & \textbf{Wynik SAW} \\ \hline
        1 & CapaCity L                 & 0.176 & 0.135 & 0.173 & 0.200 & 0.100 & \textbf{0.785} \\ \hline
        2 & Urbino 18 CNG              & 0.157 & 0.169 & 0.200 & 0.150 & 0.083 & \textbf{0.759} \\ \hline
        3 & Urbino 15 LE electric      & 0.200 & 0.200 & 0.140 & 0.100 & 0.091 & \textbf{0.731} \\ \hline
        4 & MAN Lion's Intercity LE 14 & 0.130 & 0.194 & 0.127 & 0.150 & 0.111 & \textbf{0.712} \\ \hline
        5 & Urbino 18                  & 0.152 & 0.166 & 0.127 & 0.150 & 0.111 & \textbf{0.706} \\ \hline
        6 & Urbino 12 CNG              & 0.157 & 0.123 & 0.093 & 0.150 & 0.182 & \textbf{0.705} \\ \hline
        7 & MAN Lion's City 12 E       & 0.160 & 0.135 & 0.127 & 0.200 & 0.077 & \textbf{0.699} \\ \hline
        8 & Urbino 12                  & 0.125 & 0.049 & 0.173 & 0.150 & 0.200 & \textbf{0.698} \\ \hline
        9 & Nesobus 12                 & 0.167 & 0.114 & 0.193 & 0.150 & 0.069 & \textbf{0.693} \\ \hline
        10 & SANCITY 12LF              & 0.136 & 0.123 & 0.080 & 0.150 & 0.200 & \textbf{0.689} \\ \hline
        11 & SANCITY 12LFH             & 0.169 & 0.123 & 0.160 & 0.150 & 0.062 & \textbf{0.663} \\ \hline
        12 & MAN Lion's City 12 G      & 0.137 & 0.114 & 0.060 & 0.200 & 0.100 & \textbf{0.611} \\ \hline
    \end{tabular}
    \end{adjustbox}
\end{table}

W \textbf{scenariuszu oszczędnym} zdecydowanie dominują najtańsze pojazdy -- \textit{Urbino 12} oraz \textit{SANCITY 12LF} (po 1,00 mln zł). W podstawowym modelu SAW zajmowały one odległe pozycje. Zwycięzca scenariusza bazowego, \textit{Urbino 18 CNG}, spada na siódme miejsce. Pokazuje to wyraźnie, że jego wysoka nota wynikała głównie z walorów jakościowych, a nie ekonomicznych.

W \textbf{scenariuszu technicznym} (gdzie łączna waga mocy i liczby miejsc wynosi 0,70) prym wiedzie \textit{Urbino 15 LE electric} z wynikiem 0,868. Pojazd ten dysponuje najpotężniejszym silnikiem (408 KM) i największą liczbą miejsc (65) spośród wszystkich rozpatrywanych modeli. Drugą lokatę zajmuje \textit{CapaCity L}, trzecią -- \textit{Urbino 18 CNG}. Zestawienie to nie odbiega znacząco od rankingu bazowego.

W \textbf{scenariuszu neutralnym} (wszystkie wagi równe 0,20) najlepiej wypadają autobusy o zrównoważonych parametrach. Zwycięża \textit{CapaCity L}, który łączy wysokie oceny za moc, udogodnienia dla niepełnosprawnych oraz prestiż. Warto dodać, że \textit{MAN Lion's City 12 G} zajmuje ostatnie miejsce we wszystkich trzech scenariuszach -- świadczy to o jego słabej konkurencyjności niezależnie od przyjętych preferencji. Natomiast \textit{CapaCity L} regularnie plasuje się w pierwszej trójce (poza scenariuszem oszczędnym), co potwierdza jego ogólną atrakcyjność.


\section{Średnie autobusy}

Analogiczną procedurę przeprowadzono dla segmentu autobusów średnich, stosując takie same trzy zestawy wag.

\begin{table}[H]
    \centering
    \caption{Scenariusze wag -- autobusy średnie}
    \vspace{0.2cm}
    \begin{tabular}{|l|c|c|c|c|c|}
        \hline
        \textbf{Scenariusz} & \textbf{K1: Moc} & \textbf{K2: Miejsca} & \textbf{K3: Prestiż} & \textbf{K4: Udogod.} & \textbf{K5: Koszt} \\ \hline
        Oszczędny  & 0.14 & 0.14 & 0.14 & 0.08 & 0.50 \\ \hline
        Techniczny & 0.35 & 0.35 & 0.10 & 0.15 & 0.05 \\ \hline
        Neutralny  & 0.20 & 0.20 & 0.20 & 0.20 & 0.20 \\ \hline
    \end{tabular}
\end{table}

\begin{table}[H]
    \centering
    \caption{Wyniki analizy wrażliwości -- autobusy średnie (scenariusz oszczędny)}
    \vspace{0.2cm}
    \begin{adjustbox}{max width=\textwidth}
    \begin{tabular}{|c|l|c|c|c|c|c|c|}
        \hline
        \textbf{Poz.} & \textbf{Model Autobusu} & \textbf{K1} & \textbf{K2} & \textbf{K3} & \textbf{K4} & \textbf{K5} & \textbf{Wynik SAW} \\ \hline
        1 & Crossway Pro 10,8      & 0.140 & 0.140 & 0.102 & 0.048 & 0.333 & \textbf{0.764} \\ \hline
        2 & SANCITY 9LE            & 0.082 & 0.048 & 0.070 & 0.064 & 0.500 & \textbf{0.763} \\ \hline
        3 & Crossway 10.8LE        & 0.128 & 0.086 & 0.140 & 0.064 & 0.273 & \textbf{0.691} \\ \hline
        4 & Citaro K               & 0.116 & 0.095 & 0.135 & 0.080 & 0.231 & \textbf{0.657} \\ \hline
        5 & SANCITY 10LF           & 0.096 & 0.077 & 0.092 & 0.048 & 0.333 & \textbf{0.646} \\ \hline
        6 & Urbino 8,9 LE          & 0.097 & 0.098 & 0.108 & 0.064 & 0.273 & \textbf{0.640} \\ \hline
        7 & MAN Lion's City 10 E   & 0.127 & 0.107 & 0.081 & 0.064 & 0.115 & \textbf{0.494} \\ \hline
        8 & SANCITY 10LFE          & 0.132 & 0.071 & 0.081 & 0.064 & 0.136 & \textbf{0.485} \\ \hline
        9 & Urbino 8,9 LE electric & 0.085 & 0.080 & 0.129 & 0.064 & 0.125 & \textbf{0.483} \\ \hline
        10 & Urbino 10,5 electric  & 0.132 & 0.098 & 0.065 & 0.048 & 0.115 & \textbf{0.459} \\ \hline
    \end{tabular}
    \end{adjustbox}
\end{table}

\begin{table}[H]
    \centering
    \caption{Wyniki analizy wrażliwości -- autobusy średnie (scenariusz techniczny)}
    \vspace{0.2cm}
    \begin{adjustbox}{max width=\textwidth}
    \begin{tabular}{|c|l|c|c|c|c|c|c|}
        \hline
        \textbf{Poz.} & \textbf{Model Autobusu} & \textbf{K1} & \textbf{K2} & \textbf{K3} & \textbf{K4} & \textbf{K5} & \textbf{Wynik SAW} \\ \hline
        1 & Crossway Pro 10,8      & 0.350 & 0.350 & 0.073 & 0.090 & 0.033 & \textbf{0.896} \\ \hline
        2 & Citaro K               & 0.291 & 0.238 & 0.096 & 0.150 & 0.023 & \textbf{0.798} \\ \hline
        3 & Crossway 10.8LE        & 0.321 & 0.216 & 0.100 & 0.120 & 0.027 & \textbf{0.784} \\ \hline
        4 & MAN Lion's City 10 E   & 0.317 & 0.268 & 0.058 & 0.120 & 0.012 & \textbf{0.774} \\ \hline
        5 & Urbino 10,5 electric   & 0.331 & 0.246 & 0.046 & 0.090 & 0.012 & \textbf{0.724} \\ \hline
        6 & Urbino 8,9 LE          & 0.243 & 0.246 & 0.077 & 0.120 & 0.027 & \textbf{0.713} \\ \hline
        7 & SANCITY 10LFE          & 0.331 & 0.179 & 0.058 & 0.120 & 0.014 & \textbf{0.701} \\ \hline
        8 & Urbino 8,9 LE electric & 0.212 & 0.201 & 0.092 & 0.120 & 0.013 & \textbf{0.638} \\ \hline
        9 & SANCITY 10LF           & 0.240 & 0.194 & 0.065 & 0.090 & 0.033 & \textbf{0.622} \\ \hline
        10 & SANCITY 9LE           & 0.204 & 0.119 & 0.050 & 0.120 & 0.050 & \textbf{0.543} \\ \hline
    \end{tabular}
    \end{adjustbox}
\end{table}

\begin{table}[H]
    \centering
    \caption{Wyniki analizy wrażliwości -- autobusy średnie (scenariusz neutralny)}
    \vspace{0.2cm}
    \begin{adjustbox}{max width=\textwidth}
    \begin{tabular}{|c|l|c|c|c|c|c|c|}
        \hline
        \textbf{Poz.} & \textbf{Model Autobusu} & \textbf{K1} & \textbf{K2} & \textbf{K3} & \textbf{K4} & \textbf{K5} & \textbf{Wynik SAW} \\ \hline
        1 & Crossway Pro 10,8      & 0.200 & 0.200 & 0.146 & 0.120 & 0.133 & \textbf{0.799} \\ \hline
        2 & Citaro K               & 0.166 & 0.136 & 0.192 & 0.200 & 0.092 & \textbf{0.787} \\ \hline
        3 & Crossway 10.8LE        & 0.183 & 0.123 & 0.200 & 0.160 & 0.109 & \textbf{0.776} \\ \hline
        4 & Urbino 8,9 LE          & 0.139 & 0.140 & 0.154 & 0.160 & 0.109 & \textbf{0.702} \\ \hline
        5 & MAN Lion's City 10 E   & 0.181 & 0.153 & 0.115 & 0.160 & 0.046 & \textbf{0.656} \\ \hline
        6 & SANCITY 9LE            & 0.117 & 0.068 & 0.100 & 0.160 & 0.200 & \textbf{0.645} \\ \hline
        7 & SANCITY 10LF           & 0.137 & 0.111 & 0.131 & 0.120 & 0.133 & \textbf{0.632} \\ \hline
        8 & Urbino 8,9 LE electric & 0.121 & 0.115 & 0.185 & 0.160 & 0.050 & \textbf{0.631} \\ \hline
        9 & SANCITY 10LFE          & 0.189 & 0.102 & 0.115 & 0.160 & 0.055 & \textbf{0.621} \\ \hline
        10 & Urbino 10,5 electric  & 0.189 & 0.140 & 0.092 & 0.120 & 0.046 & \textbf{0.588} \\ \hline
    \end{tabular}
    \end{adjustbox}
\end{table}

W segmencie autobusów średnich \textit{Crossway Pro 10,8} marki Iveco Bus okazuje się rozwiązaniem wyjątkowo stabilnym -- zajmuje pierwsze miejsce we wszystkich trzech scenariuszach. W scenariuszu technicznym osiąga wynik 0,896, będący najwyższą wartością spośród wszystkich badanych konfiguracji, co wynika z połączenia największej mocy silnika (360~KM) i najwyższej liczby miejsc (47) w tej klasie. W \textbf{scenariuszu oszczędnym} przewaga jest minimalna: \textit{Crossway Pro 10,8} uzyskuje wynik 0,764, wyprzedzając jedynie o 0,001 punkt \textit{SANCITY 9LE} -- najtańszy pojazd w zestawieniu (0,60~mln~zł). Oznacza to, że przy silnym priorytecie kosztowym ranking jest wrażliwy na tę różnicę i wybór decydenta może być niejednoznaczny.

\textit{Citaro K} marki Mercedes-Benz konsekwentnie plasuje się w czołowej trójce w scenariuszu technicznym (poz.~2) i neutralnym (poz.~2), potwierdzając swoją wszechstronność. \textit{SANCITY 9LE} zaś wyróżnia się wyłącznie w scenariuszu oszczędnym -- jego bardzo niska cena zakupu (0,60~mln~zł) pozwala mu wypłynąć na pozycję drugą, mimo słabszych parametrów technicznych. We wszystkich pozostałych konfiguracjach zajmuje ostatnie miejsce.

Podsumowując analizę wrażliwości obu segmentów: rekomendacja \textit{Crossway Pro 10,8} dla autobusów średnich jest odporna na zmianę preferencji decydenta, natomiast w segmencie dużych autobusów wybór lidera jest bardziej wrażliwy na strukturę wag -- w zależności od scenariusza na pozycji pierwszej pojawia się \textit{Urbino 12}, \textit{Urbino 15 LE electric} lub \textit{CapaCity L}.
% -*- TeX-master: "../main.tex" -*-
\chapter{Optymalizacja rozkładu jazdy}
Po wyborze optymalnych dla tego problemu środków transportu, należy przejść do głównego algorytmu rozwiązującego problem z wyznaczeniem rozkładu jazdy dla miasta Sławków. Algorytm jaki zastosowano do wyznaczenia opytmalnego rozkładu to \textbf{algorytm genetyczny}. Implementację tego algorytmu wykonano w środowisku \textbf{.NET} wraz z bindingiem do biblioteki \textbf{raylib}, która umożliwiła
zespołowi projektowemu zaprezentować rozwiązany problem w formie graficznej.  

\section{Algorytm genetyczny}
Algorytm genetyczny jest metodą optymalizacji inspirowaną procesami ewolucji biologicznej, w której populacja rozwiązań kandydujących ewoluuje w kierunku coraz lepszego dopasowania do zadanych kryteriów. W kontekście projektowania sieci komunikacyjnej dla miasta Sławków, każdy osobnik (rozwiązanie) reprezentuje kompletny zestaw linii autobusowych wraz z ich rozkładami jazdy. Algorytm iteracyjnie przekształca populację przy użyciu operatorów genetycznych, takich jak krzyżowanie oraz mutacja, dążąc do minimalizacji kosztów operacyjnych i maksymalizacji jakości obsługi pasażerów w ramach zdefiniowanej funkcji celu.
\subsection{Implementacja}

\subsubsection{Wagi}
Zaprojektowany model posiada następujące wagi:
\begin{itemize}
    \item \texttt{WAGA\_DLUGOSC\_TRASY}: Określa karę za zbyt długie lub nieefektywne trasy autobusowe.
    \item \texttt{WAGA\_FLOTA}: Reguluje koszt utrzymania i wielkość floty pojazdów przypisanych do linii.
    \item \texttt{KARA\_NIESPELNIONE}: Nakłada dotkliwą karę za niespełnienie wymogów funkcjonalnych projektu sieci.
    \item \texttt{WAGA\_DUPLIKATY\_DROG}: Zniechęca algorytm do tworzenia linii przebiegających przez te same odcinki dróg.
    \item \texttt{WAGA\_ZAWRACANIE}: Karze za nieoptymalne trasy wymuszające nawracanie pojazdów.
    \item \texttt{WAGA\_LICZBA\_LINII}: Kontroluje całkowitą liczbę linii autobusowych w sieci.
    \item \texttt{WAGA\_KOSZT}: Uwzględnia ekonomiczne aspekty eksploatacji całego systemu komunikacji.
    \item \texttt{SZANSA\_SWAP}: Definiuje prawdopodobieństwo wykonania operacji zamiany elementów w genotypie.
    \item \texttt{KARA\_NIEKRYCIE}: Nakłada wysoką karę za obszary, które nie zostały obsłużone przez żadną linię.
    \item \texttt{KARA\_BRAK\_POLACZ}: Karze za brak wymaganych połączeń przesiadkowych między punktami.
    \item \texttt{KARA\_SZKOLA\_CZAS}: Dotyczy niedostosowania rozkładu jazdy do godzin pracy placówek edukacyjnych.
    \item \texttt{KARA\_PRACA\_CZAS}: Dotyczy niedostosowania rozkładu jazdy do godzin dojazdów pracowników.
\end{itemize}
\subsubsection{Funkcja Fitness}
Funkcja przystosowania (\textit{fitness function}) określa jakość każdego osobnika w populacji poprzez sumowanie kar za naruszenia założeń modelu. Wartość ta jest mnożona przez -1, ponieważ algorytm dąży do minimalizacji kosztu (kary):
\begin{lstlisting}
return -(karaTrasa + karaFlota + karaObsluga + calkowityCzas * 0.01f + karaLiczbaLinii + karaKoszt + karaMaxDystans + karaNiepokryte + karaPolaczenia);
\end{lstlisting}
\subsubsection{Krok ewolucji}
Program rozpoczyna sie od stworzenia wybranej ilości linii autobusowych w tym przypadku wynosi 15 tras. Wybierani są "rodzice", na których przeprowadzane jest krzyżowanie genotypów (w przypadku tej implementacji jest to wybór losowy). Następnie na każdym z osobników populacji przeprowadzana jest mutacja zmieniająca jego genotyp, po czym przy wykorzystaniu wspomnianej wcześniej funkcji \textbf{Fitness} wybierany jest najlepszy organizm ze wszystkich możliwych.
Najlepsze genotypy są powielane, a reszta usuwana.
Dzięki temu zabiegowi otrzymywane jest następne pokolenie, które przechodzi przez taką samą procedure jak każda ubiegła epoka.
\subsubsection{Mutacja}
Proces mutacji wprowadza losowe zmiany w genotypie, umożliwiając eksplorację przestrzeni rozwiązań. Działania te obejmują trzy główne obszary:

\begin{itemize}
    \item \textbf{Dodawanie/usuwanie lini:} Dynamiczne dodawanie, usuwanie lub zamiana linii w celu zbliżenia liczby tras do wartości \texttt{CEL\_LINII}. W ten sposób decydent może zachęcić program do dobrania liczby linii bliskiej jego idelnej liczbie linii
    \item \textbf{Przebieg tras:} Modyfikacja listy przystanków (dodawanie/usuwanie) oraz optymalizacja kolejności przystanków wewnątrz linii za pomocą operacji \textit{swap}.
    \item \textbf{Rozkład jazdy:} Przesuwanie godzin odjazdu w czasie lub zmiana ich częstotliwości (dodawanie/usuwanie kursów) w ramach określonych limitów.
\end{itemize}

Intensywność tych zmian jest sterowana przez parametry \texttt{mutacja} oraz \texttt{SZANSA\_SWAP}, co zapewnia równowagę między stabilnością a poszukiwaniem lepszych wariantów sieci komunikacyjnej.

\subsubsection{Zabezpieczenie przed lokalnym ekstremum}
W celu uniknięcia przedwczesnej zbieżności algorytmu, wprowadzono mechanizmy monitorowania stagnacji populacji. Jeśli przez określoną liczbę pokoleń najlepszy osobnik nie ulega poprawie, algorytm aktywuje jedną z procedur naprawczych:

\begin{itemize}
    \item \textbf{Tryb elitarny (Nie używany):} Wyłączenie selekcji elitarnej przy jednoczesnym zwiększeniu liczebności populacji, co sprzyja szerokiej eksploracji przestrzeni rozwiązań. Jednak wyłączenie elit doprowadza do ostrego spadku
    \item \textbf{Tryb mutacji (rekomendowany):} Tryb próbuje wymusić znalesienie nowego osobnika po przez modyfikacje współczynnika mutacji, jest mniej efektywny niż elitarny ale bardziej stabilny.
    \item \textbf{Tryb automatyczny:} Poprawiona wersja trybu elitarnego, najbadziej skomplikowany tryb w naszym programie, działa w 3 trybach w przypadku braku zmiany fitness'u przez ostatni 50 generacji zwiększa populacja 4 ktotnie, po dojściu do 90 zwiąksza populacje 8 krotnie a po dojściu do 100 generacji bez poprawy przechodzi w tryb desperacji i wyłącza tymczasowo elity. Jest to tryb nie stabilny ale ma potęcjał dostarcenia lepszych wybuków niż inne tryby.
    \item \textbf{Nieskończona ilość mutacji:} Mutacje w naszej implementacji technicznie mogą trwać w nieskończoność ponieważ dzieją się w pętli, która po każdej muacji losuje liczbę od 0 do 1 i jeśli liczba jest większa od obecnego współczynnika mutacji przerywa mutowanie, pozwala na gwałtowne "wypchnięcie" osobnika z obszaru suboptymalnego.
\end{itemize}

Dzięki tym zabiegom algorytm utrzymuje zdolność do ciągłego poszukiwania globalnego optimum, zapobiegając utknięciu w lokalnych punktach krytycznych funkcji fitness.

\subsection{Wnioski post-optymalizacyjne}
Wyciągając wnioski z przeprowadzonej optymalizacji, należy zwrócić szczególną uwagę na zjawisko zbytniego usztywnienia parametrów środowiska. Jeśli model otrzyma zbyt rygorystyczne ograniczenia (np. skrajnie wysokie kary za maksymalną długość trasy lub brak tolerancji dla duplikatów dróg), może mieć trudności ze zbieżnością do akceptowalnego wyniku, co w skrajnych przypadkach skutkuje niewyznaczeniem żadnej sensownej trasy, z powodu przeważających wartości ujemnych z funkcji oceny.

Drugim kluczowym aspektem jest odpowiednie balansowanie wag przystanków. Algorytm dążąc do surowej minimalizacji kosztów operacyjnych, może wykazywać tendencję do "odcinania" przystanków o bardzo niskim popycie pasażerskim. Z matematycznego punktu widzenia jest to optymalne, jednak z punktu widzenia użyteczności publicznej prowadzi do wykluczenia komunikacyjnego peryferyjnych obszarów miasta. Z tego względu niezwykle istotne było odpowiednie dostrojenie kary za pominięcie węzłów (\texttt{KARA\_NIEKRYCIE}), aby wymusić na algorytmie włączanie ich w obieg linii komunikacyjnych, nawet kosztem nieznacznego wydłużenia czasu przejazdu.

Podsumowując, optymalizacja rozkładu za pomocą algorytmu genetycznego stanowi ciągły kompromis (trade-off) pomiędzy kosztami ponoszonymi przez przewoźnika (wielkość floty, przejechane kilometry) a komfortem pasażerów (liczba przesiadek, czas podróży). Otrzymane z programu wyniki stanowią doskonałą bazę i rekomendację decyzyjną do wdrożenia docelowego rozkładu.

\subsection{Wizualizacja rekomendacji decyzyjnych}
Poniżej zaprezentowano wizualizacje poszczególnych przebiegów linii wygenerowanych przez najlepszego osobnika (rekord) na przestrzeni ewolucji algorytmu.
\MyFigure[0.8\textwidth]{images/wizualizacja/linia_0.png}{Linia 1}{fig:Linia 1}
\MyFigure[0.8\textwidth]{images/wizualizacja/linia_1.png}{Linia 2}{fig:Linia 2}
\MyFigure[0.8\textwidth]{images/wizualizacja/linia_2.png}{Linia 3}{fig:Linia 3}
\MyFigure[0.8\textwidth]{images/wizualizacja/linia_3.png}{Linia 4}{fig:Linia 4}
\MyFigure[0.8\textwidth]{images/wizualizacja/linia_4.png}{Linia 5}{fig:Linia 5}
\MyFigure[0.8\textwidth]{images/wizualizacja/linia_5.png}{Linia 6}{fig:Linia 6}
\MyFigure[0.8\textwidth]{images/wizualizacja/linia_6.png}{Linia 7}{fig:Linia 7}
\MyFigure[0.8\textwidth]{images/wizualizacja/linia_7.png}{Linia 8}{fig:Linia 8}
\MyFigure[0.8\textwidth]{images/wizualizacja/linia_8.png}{Linia 9}{fig:Linia 9}
\MyFigure[0.8\textwidth]{images/wizualizacja/linia_9.png}{Linia 10}{fig:Linia 10}
\MyFigure[0.8\textwidth]{images/wizualizacja/linia_10.png}{Linia 11}{fig:Linia 11}
\MyFigure[0.8\textwidth]{images/wizualizacja/linia_11.png}{Linia 12}{fig:Linia 12}
\MyFigure[0.8\textwidth]{images/wizualizacja/linia_12.png}{Linia 13}{fig:Linia 13}
\MyFigure[0.8\textwidth]{images/wizualizacja/linia_13.png}{Linia 14}{fig:Linia 14}
\MyFigure[0.8\textwidth]{images/wizualizacja/linia_14.png}{Linia 15}{fig:Linia 15}

\MyFigure[1\textwidth]{images/wizualizacja/test.png}{Wszystkie linie}{fig:Wszystkie linie}

\subsection{Symulacja}
\MyFigure[0.7\textwidth]{images/wizualizacja/8000.png}{Godzina 8:00}{fig:symulacja - godzina 8}
\MyFigure[0.7\textwidth]{images/wizualizacja/1400.png}{Godzina 14:00}{fig:symulacja - godzina 14}
\MyFigure[0.7\textwidth]{images/wizualizacja/2000.png}{Godzina 20:00}{fig:symulacja - godzina 20}






% -*- TeX-master: "../main.tex" -*-
\chapter{Podsumowanie}
Zaprojektowany algorytm genetyczny pozwolił na efektywne rozwiązanie problemu optymalizacji sieci komunikacyjnej dla miasta Sławków. Dzięki zastosowaniu wielokryterialnej funkcji fitness oraz mechanizmów zabezpieczających przed utknięciem w lokalnych ekstremach, uzyskane wyniki wykazują wysoką zgodność z założonymi celami projektowymi. Implementacja w środowisku .NET z wykorzystaniem biblioteki Raylib umożliwiła nie tylko sprawne obliczenia, ale również czytelną wizualizację wypracowanego rozkładu jazdy. 

Mimo złożoności problemu optymalizacyjnego, zespół projektowy dostarczył rozwiązanie, które stanowi użyteczny model wspomagający planowanie transportu miejskiego. Wypracowane wagi oraz operatory ewolucyjne mogą służyć jako baza do dalszych prac nad udoskonalaniem systemu komunikacji publicznej w przyszłości.


% *************** Spisy pomocnicze ***************
\clearpage
\phantomsection
\addcontentsline{toc}{chapter}{Spis rysunków}
\listoffigures

\clearpage
\phantomsection
\addcontentsline{toc}{chapter}{Spis tabel}
\listoftables

\end{document}